\subsection{训练如何组织？小批次、优化器与稳定性}
\label{intro:18c-training-engineering}

训练除了定义损失函数，还要安排每次使用的数据量、更新参数的方法以及数值稳定措施。
这些设置会影响训练速度和收敛行为。

\subsubsection{Mini-batch、Batch Size 与 Epoch 是什么？}

设有 $100$ 条训练样本。如果每次取其中 $10$ 条一起估计梯度，那么 \textbf{Batch Size（批量大小）}是 $10$，这一组样本叫 \textbf{Mini-batch（小批次）}。把全部 $100$ 条训练样本完整处理一遍，称为一个 \textbf{Epoch（训练轮次）}；本例中每轮大约需要 $100/10=10$ 次参数更新。

\textbf{优化器}决定怎样利用梯度调整参数。普通梯度下降按学习率和梯度更新权重；\textbf{Adam} 会进一步使用梯度的一阶与二阶矩估计，自适应地调节更新幅度，但不保证在所有任务上表现最好。

若梯度特别大，\textbf{梯度裁剪}可以在更新前限制其幅度。例如梯度向量的长度超过设定阈值时，可以将其整体按比例缩小。裁剪不能取代合理的学习率和数据检查。

\noindent\textbf{先记主线：}SGD 用当前批次的梯度和学习率直接更新参数；出现不稳定时，先检查数据、梯度大小与学习率，而不是一次背完所有优化器。动量、Adam、AdamW 与 Warmup 的独立手算已集中到\hyperref[app:optimizer-details]{附录 A 的提高选读}。\par\medskip

\Needspace{24\baselineskip}
\subsubsection{训练不稳定时先诊断什么？}

\textbf{梯度消失}先排查初始化、激活与残差路径；
\textbf{梯度爆炸}检查学习率并考虑裁剪；
\textbf{更新震荡}核对学习率、动量与 Adam；
\textbf{样本过多}则通过 Mini-batch 和硬件并行降低单次更新成本。
这些操作分别处理数值传播、优化步长或吞吐问题，
不能互相替代。

\textbf{超参数}是训练前或训练过程中设定的控制量，如学习率、正则化系数、批次大小；\textbf{模型参数}则是通过优化直接学习的权重。

\begin{example}
\noindent\emph{一次实际的梯度裁剪}\par

设当前参数梯度为 $(6,8)$，其 L2 长度是 $10$。如果阈值设为 $5$，整体乘以 $5/10=0.5$，得到裁剪后的梯度 $(3,4)$；方向不变，长度变为 $5$。用学习率 $0.1$ 更新时，两个参数分别减去 $0.3$ 和 $0.4$。裁剪的对象是梯度，不是原始训练数据。
\end{example}

\chaptercheck{16-loss-gradient}{梯度下降的更新方向与损失梯度是什么关系？}
{三点损失从 $(w,b)=(1,0)$ 出发，梯度是多少？取 $\eta=0.3$，更新后的损失是多少？}
{学习率取零会怎样？若梯度为 $(6,8)$，按 L2 阈值 $5$ 裁剪后变成多少？}

\begin{conclusion}
我们已用三点数据算出损失、梯度和同步更新；
批次、学习率调度与优化器进一步决定训练怎样执行。
接下来研究：若函数只是一条直线，怎样表示弯曲的关系？
\end{conclusion}
